%!TEX TS-program = xelatex

% Данный шаблон подготовлен для курса LaTeX в РАНХиГС
% на основе шаблона 
% Данилы Фёдоровых (danil@fedorovykh.ru),
%  который использовал его в курсе 
% <<Документы и презентации в \LaTeX>> НИУ ВШЭ
% Исходная версия шаблона --- 
  % https://www.writelatex.com/coursera/latex/5.1
\documentclass[c, dvipsnames]{beamer}  % [t], [c], или [b] --- вертикальное 
%\documentclass[handout, dvipsnames, c]{beamer} % Раздаточный материал (на слайдах всё сразу)
% !TEX root = main_file.tex

%%%%%%%%%% Програмный код %%%%%%%%%%
% \usepackage{minted}
% Включает подсветку команд в программах!
% Нужно, чтобы на компе стоял питон, надо поставить пакет Pygments, в котором он сделан, через pip.

% Для Windows: Жмём win+r, вводим cmd, жмём enter. Открывается консоль.
% Прописываем pip install Pygments
% Заходим в настройки texmaker и там прописываем в PdfLatex или XelaTeX:
% pdflatex -shell-escape -synctex=1 -interaction=nonstopmode %.tex

% Для Linux: Открываем консоль. Убеждаемся, что у вас установлен pip командой pip --version
% Если он не установлен, ставим его: sudo apt-get install python-pip
% Ставим пакет sudo pip install Pygments

% Для Mac: Всё то же самое, что на Linux, но через brew.

% После всего этого вы должны почувствовать себя тру-программистами!
% Документация по пакету хорошая. Сам читал, погуглите!

%%%%%%%%%% Математика %%%%%%%%%%
\usepackage{amsmath,amsfonts,amssymb,amsthm,mathtools}
% Показывать номера только у тех формул, на которые есть \eqref{} в тексте.
%\mathtoolsset{showonlyrefs=true}
%\usepackage{leqno} % Нумерация формул слева

%%%%%%%%%% Шрифты %%%%%%%%%%
\usepackage[english, russian]{babel}   % выбор языка для документа
\usepackage[utf8]{inputenc}            % выбор utf8 кодировки
\usepackage[X2,T2A]{fontenc}           % ещё немного кодировки

\usepackage{fontspec}                  % пакет для подгрузки шрифтов
\setmainfont{Times New Roman}          % задаёт основной шрифт документа

\usepackage{unicode-math}              % пакет для установки математического шрифта
\setmathfont{Asana-Math.otf}           % шрифт для математики

% Конкретный символ из конкретного шрифта
% \setmathfont[range=\int]{Neo Euler}

%%%%%%%%%% Работа с картинками %%%%%%%%%%
\usepackage{graphicx}                  % Для вставки рисунков
\usepackage{graphics}
\graphicspath{{images/}{pictures/}}    % можно указать папки с картинками
\usepackage{wrapfig}                   % Обтекание рисунков и таблиц текстом


%%%%%%%%%% Работа с таблицами %%%%%%%%%%
\usepackage{tabularx}            % новые типы колонок
\usepackage{tabulary}            % и ещё новые типы колонок
\usepackage{array,delarray}      % Дополнительная работа с таблицами
\usepackage{longtable}           % Длинные таблицы
\usepackage{multirow}            % Слияние строк в таблице
\usepackage{float}               % возможность позиционировать объекты в нужном месте
\usepackage{booktabs}            % таблицы как в книгах
% Заповеди из документации к booktabs:
% 1. Будь проще! Глазам должно быть комфортно
% 2. Не используйте вертикальные линни
% 3. Не используйте двойные линии. Как правило, достаточно трёх горизонтальных линий
% 4. Единицы измерения - в шапку таблицы
% 5. Не сокращайте .1 вместо 0.1
% 6. Повторяющееся значение повторяйте, а не говорите "то же"
% 7. Есть сомнения? Выравнивай по левому краю!

%  вычисляемые колонки по tabularx
\newcolumntype{C}{>{\centering\arraybackslash}X}
\newcolumntype{L}{>{\raggedright\arraybackslash}X}
\newcolumntype{Y}{>{\arraybackslash}X}
\newcolumntype{Z}{>{\centering\arraybackslash}X}


%%%%%%%%%% Графика и рисование %%%%%%%%%%
\usepackage{tikz, pgfplots}    % язык для рисования графики из latex'a

%%%%%%%%%% Гиперссылки %%%%%%%%%%
\usepackage{xcolor}            % разные цвета

\usepackage{hyperref}
\hypersetup{
    unicode=true,           % позволяет использовать юникодные символы
    colorlinks=true,       	% true - цветные ссылки, false - ссылки в рамках
    urlcolor =blue,         % цвет ссылки на url
    linkcolor=black,        % внутренние ссылки
    citecolor=black,        % на библиографию
	breaklinks              % если ссылка не умещается в одну строку, разбивать ли ее на две части?
}

%%%%%%%%%% Другие приятные пакеты %%%%%%%%%%
\usepackage{multicol}       % несколько колонок
\usepackage{verbatim}       % для многострочных комментариев
\usepackage{enumitem}       % дополнительные плюшки для списков
%  например \begin{enumerate}[resume] позволяет продолжить нумерацию в новом списке

\usepackage{todonotes}      % для вставки в документ заметок о том, что  осталось сделать
% \todo{Здесь надо коэффициенты исправить}
% \missingfigure{Здесь будет Последний день Помпеи}
% \listoftodos --- печатает все поставленные \todo'шки

%%%%%%%%%%%%%%%%%%%%%%%%%%%%%%%%%%%%%%%%%%%
%%%%%%%%%% ГОСТОВСКИЕ ПРИБАМБАСЫ %%%%%%%%%%
%%%%%%%%%%%%%%%%%%%%%%%%%%%%%%%%%%%%%%%%%%%

% размер листа бумаги
\usepackage[paper=a4paper,top=15mm, bottom=15mm,left=35mm,right=10mm,includehead]{geometry}

% всякие разные расстояния
\usepackage{setspace}
\setstretch{1.33}              % Полуторный межстрочный интервал
\setlength{\parindent}{1.5em}  % Красная строка.

\righthyphenmin=2    % Разрешение переноса двух и более символов
\widowpenalty=10000  % Наказание за вдовствующую строку (одна строка абзаца на этой странице, остальное --- на следующей)
%\clubpenalty=10000  % Наказание за сиротствующую строку (омерзительно висящая одинокая строка в начале страницы)
\tolerance=1000      % Ещё какое-то наказание.

% Нумерация страниц сверху по центру
\usepackage{fancyhdr}
\pagestyle{fancy}
\fancyhead{ } % clear all fields
\fancyfoot{ } % clear all fields
\fancyhead[C]{\thepage}
% Чтобы не прорисовывалась черта!
\renewcommand{\headrulewidth}{0pt}

% Нумерация страниц с надписью "Глава"
\usepackage{etoolbox}
\patchcmd{\chapter}{\thispagestyle{plain}}{\thispagestyle{fancy}}{}{}

% Заголовки по левому краю
% опция identfirst устанавливает отступ в первом абзаце
\usepackage[indentfirst]{titlesec}{\raggedleft}

% В Linux этот пакет для заголовков. Исправляет это следующий непонятный кусок кода:
\makeatletter
\patchcmd{\ttlh@hang}{\parindent\z@}{\parindent\z@\leavevmode}{}{}
\patchcmd{\ttlh@hang}{\noindent}{}{}{}
\makeatother

% Редактирования Глав и названий
\titleformat{\chapter}
      {\normalfont\large\bfseries}
      {\thechapter }{0.5 em}{}

% Редактирование ненумеруемых глав chapter* (Введение и тп)
\titleformat{name=\chapter,numberless}
{\centering\normalfont\bfseries\large}{}{0.25em}{\normalfont}

% Убирает чеканутые отступы вверху страницы
\titlespacing{\chapter}{0pt}{-\baselineskip}{\baselineskip}

% Более низкие уровни (подзаголовки)
\titleformat{\section}{\bfseries}{\thesection}{0.5 em}{}
\titleformat{\subsection}{\bfseries}{\thesubsection}{0.5 em}{}

\titlespacing*{\section}{0 pt}{\baselineskip}{\baselineskip}
\titlespacing*{\subsection}{0 pt}{\baselineskip}{\baselineskip}

% Содержание, команды ниже изменяют отступы и рисуют точечки!
\usepackage{titletoc}

\titlecontents{chapter}
             [1em] %
             {\normalsize}
             {\contentslabel{1 em}}
             {\hspace{-1 em}}
             {\normalsize\titlerule*[10pt]{.}\contentspage}

\titlecontents{section}
              [3 em] %
              {\normalsize}
              {\contentslabel{1.75 em}}
              {\hspace{-1.75 em}}
              {\normalsize\titlerule*[10pt]{.}\contentspage}

\titlecontents{subsection}
              [6 em] %
              {\normalsize}
              {\contentslabel{3 em}}
              {\hspace{-3 em}}
              {\normalsize\titlerule*[10pt]{.}\contentspage}

% Правильные подписи под таблицей и рисунком
% Документация к пакету на русском языке!
\usepackage[tableposition=top, singlelinecheck=false]{caption}
\usepackage{subcaption}

   \DeclareCaptionStyle{base}%
		[justification=centering,indention=0pt]{}
   \DeclareCaptionLabelFormat{gostfigure}{Рисунок #2}
   \DeclareCaptionLabelFormat{gosttable}{Таблица #2}

   \DeclareCaptionLabelSeparator{gost}{~---~}
   \captionsetup{labelsep=gost}

   \DeclareCaptionStyle{fig01}%
           [margin=5mm,justification=centering]%
           {margin={3em,3em}}
   \captionsetup*[figure]{style=fig01,labelsep=gost,labelformat=gostfigure,format=hang}

   \DeclareCaptionStyle{tab01}%
           [margin=5mm,justification=centering]%
           {margin={3em,3em}}
   \captionsetup*[table]{style=tab01,labelsep=gost,labelformat=gosttable,format=hang}

% межстрочный отступ в таблице
 \renewcommand{\arraystretch}{1.2}

% многостраничные таблицы под РОССИЙСКИЙ СТАНДАРТ
% ВНИМАНИЕ! Обязательно после пакета caption!
\usepackage{fr-longtable}

%Более гибкие спсики
\usepackage{enumitem}

% сообщаем окружению о том, что существует такая штука, как нумерация русскими буквами
\makeatletter
\AddEnumerateCounter{\asbuk}{\russian@alph}{щ}
\makeatother

% ГОСТОВСКИЕ СПИСКИ

% Первый тип списков. Большая буква.
\newlist{Enumerate}{enumerate}{1}

\setlist[Enumerate,1]{labelsep=0.5em,leftmargin=1.25em,labelwidth=1.25em,
parsep=0em,itemsep=0em,topsep=0ex, before={\parskip=-1em},label=\arabic{Enumeratei}.}

% Второй тип списков. Маленькая буква.
\setlist[enumerate]{label=\arabic{enumi}),parsep=0em,itemsep=0em,topsep=0.75ex, before={\parskip=-1em}}

% Третий тип списков. Два уровня.
\newlist{twoenumerate}{enumerate}{2}
\setlist[twoenumerate,1]{itemsep=0mm,parsep=0em,topsep=0.75ex,, before={\parskip=-1em},label=\asbuk{twoenumeratei})}
\setlist[twoenumerate,2]{leftmargin=1.3em,itemsep=0mm,parsep=0em,topsep=0ex, before={\parskip=-1em},label=\arabic{twoenumerateii})}

% Четвёртый тип списков. Список с тире.
\setlist[itemize]{label=--,parsep=0em,itemsep=0em,topsep=0ex, before={\parskip=-1em},after={\parskip=-1em}}

% WARNING WARNING WARNIN!
% Если в списке предложения, то должна по госту стоять точка после цифры => команда Enumerate! Если идет перечень маленьких фактов, не обособляемых предложений то после цифры идет скобка ")" => команда enumerate! Если перечень при этом ещё и двууровневый, то twoenumerate.


%%%%%%%%%% Список литературы %%%%%%%%%%

%\usepackage[
%backend=biber, %подключение пакета biber (тоже нужен)
%bibstyle=gost-numeric, %подключение одного из четырех главных стилей biblatex-gost
%sorting=ntvy, %тип сортировки в библиографии
%]{biblatex}
\usepackage[backend=biber,style=gost-numeric, maxbibnames=9,maxcitenames=2,uniquelist=false, babel=other]{biblatex}

% Справка по 4 главным стилям для ленивых:
% gost-inline  ссылки внутри теста в круглых скобках
% gost-footnote подстрочные ссылки
% gost-numeric затекстовые ссылки
% gost-authoryear тоже затекстовые ссылки, но немного другие

% Подробнее смотри страницу 4 документации. Она на русском.

% Ещё немного настроек
\DeclareFieldFormat{postnote}{#1} %убирает с. и p.
\renewcommand*{\mkgostheading}[1]{#1} % только лишь убираем курсив с авторов

\addbibresource{bib.bib} 

% Этот кусок кода выносит русские источники на первое место. Костыль описали авторы пакета в руководстве к нему. Подробнее смотри:
% https://github.com/odomanov/biblatex-gost/wiki/Как-сделать%2C-чтобы-русскоязычные-источники-предшествовали-остальным
\DeclareSourcemap{
  \maps[datatype=bibtex]{
    \map{
      \step[fieldsource=langid, match=russian, final]
      \step[fieldset=presort, fieldvalue={a}]
    }
    \map{
      \step[fieldsource=langid, notmatch=russian, final]
      \step[fieldset=presort, fieldvalue={z}]
    }
  }
}

\DefineBibliographyStrings{english}{%
pages = {P\adddot},
number = {№},
}

\title[Эмпирический анализ моделей эк. роста ]{Эмпирический анализ выводов \\ моделей экономического роста}
\subtitle{Защита выпускной квалификационной работы}


\author[Александр Тишин]{Александр Тишин \\ \smallskip \scriptsize ЭО-13-02 \\ \smallskip \scriptsize \href{mailto:phenyard@gmail.com}{\nolinkurl{phenyard@gmail.com} }}

\superviser{к.э.н. Перевышин Ю.Н.}

%\author[Имя автора]{Имя автора \\ \smallskip \scriptsize \href{mailto:author@ranepa.ru}{author@ranepa.ru} \\ \smallskip  \href{http://ranepa.ru}{http://ranepa.ru} }

\institute[РАНХиГС]{ \uppercase{
  Российская Академия Народного Хозяйства и  \\ Государственной Службы при Президенте Российской Федерации}}
\date{23 июня 2017 г.}


\titlegraphic{\includegraphics[scale=0.5]{logo1}}
\titlegraphicii{\includegraphics[scale=0.5]{logo2}}

\begin{document}

\frame[plain]{\titlepage}	% Титульный слайд


\begin{frame}[c]{Актуальность исследования} 
\begin{itemize}
\item  Зуб даю, актуально! 
  \item  Ничего нет актуальнее!!!
  \item  Актуальность --- зашкаливает
\item  Актуальная актуальность с актуальным актуалитетом.
\end{itemize}
\end{frame}

%
%
%\begin{frame}[c]
%\frametitle{Анализ предметной области}
%{ \small   % вместо small можно поставить scriptsize чтобы влезло
%	\begin{table}[]
%		\centering
%		\resizebox{\textwidth}{!}{ 
  %			\begin{tabular}{|p{2.2cm}|p{1.8cm}|p{3.5cm}|p{7cm}|}
  %				\hline\rowcolor{backgr}
  %				\textcolor{white}{Авторы} & \textcolor{white}{Выборка, период}  & \textcolor{white}{Метод исследования}& \textcolor{white}{Результат} \\			
  %				\hline
  %				(Kuper, 2003)  & США, 998-2008  &  Коинтеграци и VECM & Получились значимые результаты с интересной интерпретацией.\\
  %				\hline
  %		\end{tabular} }
%	\end{table}
%}
%\end{frame}
%
%
%\begin{frame}[shrink=3]
%	\frametitle{Цели и задачи}
%	\begin{block}{Цель:}
%	\begin{itemize}
%		\item Единственная цель
%	\end{itemize}
%		
%	\end{block}
%
%	 	\begin{block}{Задачи:}
%			\begin{enumerate}
%	\item Найти красивый шаблон
%	\item Заручиться поддержкой научного руководителя 
%	\item Ходить на научные семинары
%	\item ?????
%	\item Profit
%	 \end{enumerate}	
%	\end{block}
%\end{frame}
%
%
%
%\section{Неоклассические модели экономического роста}
%\subsection{Страны, пространственная выборка}
%
%\begin{frame}[shrink=5]
%\frametitle{\insertsection} 
%\framesubtitle{\insertsubsection}
%Модель (Mankiw N. G., Romer D., Weil D. N., 1992): 
  %\begin{equation} 
%\log(y_i)=a+\frac{\alpha}{1-\alpha}\log((s_k)_i)-\frac{\alpha}{1-\alpha}\log(n_i+g+\delta)+\varepsilon_i
%\end{equation}
%где: 
  %\begin{itemize}
%\item $y$ "--- ВВП на душу населения по ППС  в долларах США в 2013 гг.
%\item $s_k$ "--- средняя доля инвестиций в  ВВП (Gross capital formation), так как $I=S$
  %\item $n$ "--- средний темп прироста  численности населения  ($g+\delta=5\%$)
%\end{itemize}
%Рассматриваются три группы стран:
  %\begin{itemize}
%\item Высокий доход, не нефтедобывающие~---~68 стран 
%\item Средний доход --- 67 страны
%\item Высокий доход, входящие в ОЭСР --- 22 страны
%\end{itemize}
%Отличия:
  %\begin{itemize}
%	\item Непересекающийся с MRW временной интервал с 1990 по 2013 гг.
%
%\end{itemize}
%\end{frame}
%
%\begin{frame}
% 			\frametitle{\insertsection}
%	\framesubtitle{\insertsubsection}
%	\vspace{-1ex}
%	\begin{itemize}
%		\item \alert{ Значимые коэффициенты }для первой и второй групп стран
%		\item \alert{Предполагаемые теорией знаки:} отрицательный перед темпами прироста населения и положительный перед долей инвестиций 
%		\item Достаточно высокий $R^2_{adj}$
  %	\end{itemize}
%	\vspace{-2ex}
%	
%\begin{table}[!htbp] \centering 
%\resizebox{0.9\textwidth}{!}{ 
  %\begin{tabular}{@{\extracolsep{5pt}}lD{.}{.}{-2} D{.}{.}{-2} D{.}{.}{-2} } 
  %\\[-1.8ex]\hline 
  %\hline \\[-1.8ex] 
  %\\[-1.8ex] & \multicolumn{3}{c}{Зависимая переменная: ВВП ($y$)} \\ 
  % & \multicolumn{1}{c}{Первая группа} & \multicolumn{1}{c}{Вторая группа} & \multicolumn{1}{c}{Третья группа} \\ 
  %\hline \\[-1.8ex] 
  % $\log(n+g+\delta)$ & -7.68^{***} & -6.60^{***} & -0.96 \\ 
  %  & (0.82) & (0.87) & (1.07) \\ 
  %  $\log(s_k)$ & 2.04^{***} & 1.53^{***} & 1.20 \\ 
  %  & (0.48) & (0.57) & (1.15) \\ 
  %  Constant & 17.28^{***} & 17.01^{***} & 8.58^{**} \\ 
  %  & (2.31) & (2.55) & (3.62) \\ 
  % \textit{N} & \multicolumn{1}{c}{86} & \multicolumn{1}{c}{67} & \multicolumn{1}{c}{22} \\ 
  %
  %Adjusted R$^{2}$ & \multicolumn{1}{c}{0.58} & \multicolumn{1}{c}{0.51} & \multicolumn{1}{c}{-0.02} \\ 
  %
  %\hline 
  %\hline \\[-1.8ex] 
  %
  %\end{tabular} }
%\end{table} 
%\end{frame}
%
%
%
%\subsection{Страны, панель}
%
%
%
%\begin{frame}[shrink=12]
%\frametitle{\insertsection} 
%\framesubtitle{\insertsubsection}
%Модель панельной регрессии с фиксированными временными эффектами: 
  %\begin{equation} 
%\log(y_{i,t})=a_t+\frac{\alpha}{1-\alpha}\log((s_k)_{i,t})-\frac{\alpha}{1-\alpha}\log(n_{i,t}+g+\delta)+\varepsilon_{i,t}
%\end{equation}
%где 
%\begin{itemize}
%\item $y$ "--- средний ВВП на душу населения по ППС  в долларах США в постоянных ценах  за шесть лет
%\item $s_k$ "--- средняя доля инвестиций в  ВВП  за шесть лет
%\item $n$ "--- средний темп прироста  численности населения  за шесть лет
%\end{itemize}
%Рассматриваются три группы стран:
%\begin{itemize}
%\item Высокий доход, не нефтедобывающие~---~68 стран 
%\item Средний доход --- 67 страны
%\item Высокий доход, входящие в ОЭСР --- 22 страны
%\end{itemize}
%Построение панелей:
%\begin{itemize}
%	\item Данные с 1990 по 2013 гг. усреднялись по 6 лет $\Rightarrow$ 4 временных промежутка 
%\end{itemize}
%
%\end{frame}
%
%\begin{frame}
% 			\frametitle{\insertsection}
%	\framesubtitle{\insertsubsection}
%	\begin{itemize}
%		\item Результаты \alert{ совпадают с пространственным анализом}: статистически значимые коэффициенты и предполагаемые теорией знаки
%	\end{itemize}
%	\vspace{-2ex}
%\begin{table}[!htbp] \centering 
%\resizebox{\textwidth}{!}{ 
%\begin{tabular}{@{\extracolsep{5pt}}lD{.}{.}{-2} D{.}{.}{-2} D{.}{.}{-2} } 
%\\[-1.8ex]\hline 
%\hline \\[-1.8ex] 
%\\[-1.8ex] & \multicolumn{3}{c}{Зависимая переменная: ВВП ($y$)} \\ 
% & \multicolumn{1}{c}{Первая группа} & \multicolumn{1}{c}{Вторая группа} & \multicolumn{1}{c}{Третья группа} \\ 
%\hline \\[-1.8ex] 
% $\log(n+g+\delta)$ & -5.52^{***} & -5.36^{***} & -0.79^{*} \\ 
%  & (0.42) & (0.45) & (0.41) \\ 
%  $\log(s_k)$ & 1.43^{***} & 1.08^{***} & 1.05^{**} \\ 
%  & (0.22) & (0.26) & (0.43) \\ 
% \textit{N} & \multicolumn{1}{c}{359} & \multicolumn{1}{c}{277} & \multicolumn{1}{c}{88} \\ 
%Adjusted R$^{2}$ & \multicolumn{1}{c}{0.38} & \multicolumn{1}{c}{0.37} & \multicolumn{1}{c}{0.02} \\ 
%\hline 
%\hline \\[-1.8ex] 
%\end{tabular} }
%\end{table}
%\end{frame}
%
%\section{Выводы по неоклассическим моделям}
%\begin{frame}
%	\frametitle{\insertsection}
%	Страны:
%	\begin{itemize}
%		\item Норма сбережений физического и человеческого капитала и темп прироста населения являются факторами объясняющими различия в уровне подушевого дохода 
%		\item Модель Солоу объясняет $60\%-70\%$ вариации подушевого дохода для первой и второй групп стран. Добавление человеческого капитала улучшает выводы (\textit{уменьшаются коэффициенты})
%	\end{itemize}	
%	Регионы 
%	\begin{itemize}
%		\item Свидетельств  того, что различия подушевого дохода регионов объясняются за счет факторов неоклассических моделей экономического роста не выявлено
%		\item Возможные объяснения: нарушены предпосылки моделей, например, закрытая экономика 
%	\end{itemize}
%\end{frame}
% 
%\section{Конвергенция}
%
%
%\begin{frame}
%\frametitle{\insertsection} 
% 
%\begin{block}{Безусловная $\beta$-конвергенция }
% \begin{equation}
%\dfrac{1}{T}\log(\dfrac{y_{i,T}}{y_{i,0}})=\alpha-\beta \log(y_{i,0})+\varepsilon_{i}
% \end{equation} 
% Бедные страны растут быстрее богатых, и разница в уровнях подушевого дохода постепенно снижается независимо от  характеристик экономики
% \end{block}
% \begin{block}{Условная $\beta$-конвергенция}
%  \begin{equation}
%\dfrac{1}{T}\log(\dfrac{y_{i,T}}{y_{i,0}})=\alpha-\beta \log(y_{i,0})+\gamma X+ \varepsilon_{i}
% \end{equation} 
% Бедные страны растут быстрее богатых при условиях одинаковых характеристик экономики 
%  \end{block}
%  
%  $y_i$ --- подушевой доход для $i$ субъекта,  $0$ --- в начальный момент времени и $T$ --- в конечный момент времени
%
%
%\end{frame}
%
%
%\subsection{Страны, пространственная выборка, безусловная конвергенция}
%
%\frame<beamer:0>{
%% Можно и вернуть обратно
% 			\frametitle{\insertsection}
%	\framesubtitle{\insertsubsection}
%		\begin{itemize}
%		\item \alert{Не отвергается гипотеза о  безусловной конвергенции} для первой и второй групп: статистически значимые отрицательные коэффициенты 
%        
%	\end{itemize}
%\begin{table}[!htbp] \centering 
%\resizebox{\textwidth}{!}{ 
%\begin{tabular}{@{\extracolsep{5pt}}lD{.}{.}{-2} D{.}{.}{-2} D{.}{.}{-2} } 
%\\[-1.8ex]\hline 
%\hline \\[-1.8ex] 
%\\[-1.8ex] & \multicolumn{3}{c}{Зависимая переменная: средний темп роста ВВП} \\
% & \multicolumn{1}{c}{Первая группа} & \multicolumn{1}{c}{Вторая группа} & \multicolumn{1}{c}{Третья группа} \\ 
%\hline \\[-1.8ex] 
% $\log(y_{1990})$ & -0.14^{***} & -0.15^{***} & -0.22 \\ 
%  & (0.04) & (0.04) & (0.20) \\ 
%  Constant & 1.76^{***} & 1.89^{***} & 2.43 \\ 
%  & (0.31) & (0.37) & (2.11) \\ 
% \textit{N} & \multicolumn{1}{c}{85} & \multicolumn{1}{c}{69} & \multicolumn{1}{c}{22} \\ 
%Adjusted R$^{2}$ & \multicolumn{1}{c}{0.14} & \multicolumn{1}{c}{0.15} & \multicolumn{1}{c}{0.01} \\ 
%\hline 
%\hline \\[-1.8ex] 
%\end{tabular} }
%\end{table} 
%}
%
%\subsection{Страны, панель, безусловная конвергенция}
%
%\begin{frame}
% 			\frametitle{\insertsection}
%	\framesubtitle{\insertsubsection}
%	\begin{itemize}
% 
%		\item Статистически значимые и отрицательные  коэффициенты для всех групп $\Rightarrow$ \alert{безусловная конвергенция}
%		\item В теории из безусловной конвергенции следует условная, подробнее в раздаточных материалах 
%			\item Лучше пространственной регрессии
%	\end{itemize}
%	\vspace{-3ex}
%\begin{table}[!htbp] \centering 
%	\resizebox{\textwidth}{!}{ 
%\begin{tabular}{@{\extracolsep{5pt}}lD{.}{.}{-2} D{.}{.}{-2} D{.}{.}{-2} } 
%\\[-1.8ex]\hline 
%\hline \\[-1.8ex] 
%\\[-1.8ex] & \multicolumn{3}{c}{Зависимая переменная: средний темп роста ВВП} \\
% & \multicolumn{1}{c}{Первая группа} & \multicolumn{1}{c}{Вторая группа} & \multicolumn{1}{c}{Третья группа} \\ 
%\hline \\[-1.8ex] 
% $\log(y_{0})$ & -0.10^{***} & -0.12^{***} & -0.51^{**} \\ 
%  & (0.03) & (0.04) & (0.23) \\ 
% \textit{N} & \multicolumn{1}{c}{306} & \multicolumn{1}{c}{242} & \multicolumn{1}{c}{71} \\ 
%Adjusted R$^{2}$ & \multicolumn{1}{c}{0.02} & \multicolumn{1}{c}{0.02} & \multicolumn{1}{c}{0.01} \\ 
%\hline 
%\hline \\[-1.8ex] 
%\end{tabular} }
%\end{table}
%\end{frame}
%
%\section{Выводы по конвергенции}
%\begin{frame}
%\frametitle{\insertsection}
% 
%\begin{itemize}
%	\item Конвергенция по странам не отвергается для всех трех групп (панель)
%	\item Конвергенция по регионам не отвергается  для обеих групп 
%	\item Таким образом, как на страновом уровне, так и на региональном есть <<конвергенционные клубы>>, подушевой доход их <<участников>> имеет тенденцию к сближению 
%\end{itemize} 
%\end{frame}
% 
%
%\begin{frame}[c, plain]
%\begin{center}
%
%{\LARGE Спасибо за внимание}
%
%\bigskip
%
%{\Large \inserttitle}
%
%\bigskip
%
%{\insertauthor} 
%
%\bigskip
%
%\bigskip\bigskip
%
%{\large \insertdate}
%\end{center}
%\end{frame}
%
%\section{Приложения}
%\begin{frame}[plain,c]
%\begin{center}
%\Huge{\textbf{
%	ПРИЛОЖЕНИЯ}}
%\end{center}
%	
%\end{frame}
%
%
%\begin{frame}
%	\frametitle{\insertsection}
%	\framesubtitle{Группы регионов }
%\begin{table}[ht]
%	\centering
%	\resizebox{\textwidth}{!}{ 
%		\begin{tabular}{|lr|lr|lr|}
% \hline
% регион & N & регион & N & регион & N \\
% \hline
% Белгородская область & 1 & Брянская область & 2 & Республика Марий Эл & 2 \\ Липецкая область & 1 & Владимирская область & 2 & Республика Мордовия & 2 \\ Московская область & 1 & Воронежская область & 2 & Удмуртская Республика & 2 \\ Республика Коми & 1 & Ивановская область & 2 & Чувашская Республика & 2 \\ Вологодская область & 1 & Калужская область & 2 & Кировская область & 2 \\ Ленинградская область & 1 & Костромская область & 2 & Нижегородская область & 2 \\ Новгородская область & 1 & Курская область & 2 & Пензенская область & 2 \\ г.Санкт-Петербург & 1 & Орловская область & 2 & Саратовская область & 2 \\ Краснодарский край & 1 & Рязанская область & 2 & Ульяновская область & 2 \\ Астраханская область & 1 & Смоленская область & 2 & Курганская область & 2 \\ Республика Башкортостан & 1 & Тамбовская область & 2 & Республика Алтай & 2 \\ Республика Татарстан & 1 & Тверская область & 2 & Республика Бурятия & 2 \\ Пермский край & 1 & Тульская область & 2 & Республика Тыва & 2 \\ Оренбургская область & 1 & Ярославская область & 2 & Республика Хакасия & 2 \\ Самарская область & 1 & Республика Карелия & 2 & Алтайский край & 2 \\ Свердловская область & 1 & Калининградская область & 2 & Новосибирская область & 2 \\ Челябинская область & 1 & Мурманская область & 2 & Омская область & 2 \\ Иркутская область & 1 & Псковская область & 2 & Приморский край & 2 \\ Кемеровская область & 1 & Республика Адыгея & 2 & Хабаровский край & 2 \\ Томская область & 1 & Республика Калмыкия & 2 & Амурская область & 2 \\ Республика Саха (Якутия) & 1 & Волгоградская область & 2 & Магаданская область & 2 \\ Чукотский автономный округ & 1 & Ростовская область & 2 & Еврейская автономная область & 2 \\
% \hline
%	\end{tabular} }
%\end{table}
%
%\end{frame}
%
%
%
%\begin{frame}
%\frametitle{\insertsection}
%	\framesubtitle{Оценки $\alpha$ и $\beta$ }
%	\begin{table}[ht]
%\centering 
%	\resizebox{\textwidth}{!}{ 
%\begin{tabular}{|l|c|c|c|c|c|c|}
%  \hline
% Модель & \multicolumn{2}{c|}{Первая группа} & \multicolumn{2}{c|}{Вторая группа} & \multicolumn{2}{c|}{Третья группа} \\ 
%Солоу & GDP PPP & GDP PC & GDP PPP & GDP PC & GDP PPP & GDP PC \\ 
%  \hline
%Оценка $\alpha$ & 0.75 & 0.78 & 0.71 & 0.76 & 0.27 & 0.52 \\ 
%  Левая граница & 0.70 & 0.73 & 0.64 & 0.68 & -2.40 & -1.58 \\ 
%  Правая граница & 0.79 & 0.82 & 0.76 & 0.81 & 0.59 & 0.73 \\ 
%   \hline
%\end{tabular} }
%\end{table}
%
%\begin{table}[ht]
%\centering
%	\resizebox{\textwidth}{!}{ 
%\begin{tabular}{|l|c|c|c|c|c|c|}
%  \hline
%Модель Солоу & \multicolumn{2}{c|}{Первая группа} & \multicolumn{2}{c|}{Вторая группа} & \multicolumn{2}{c|}{Третья группа} \\ 
%с чел. капиталом & GDP PPP & GDP PC & GDP PPP & GDP PC & GDP PPP & GDP PC \\ 
%  \hline
%Оценка $\alpha$ & 0.46 & 0.45 & 0.37 & 0.34 & -0.14 & 0.19 \\ 
%  Левая граница & 0.37 & 0.33 & 0.22 & 0.12 & 2.47 & 1.78 \\ 
%  Правая граница & 0.51 & 0.51 & 0.45 & 0.45 & 0.43 & 0.56 \\ 
%   \hline
%\end{tabular} }
%\end{table}
%
%\begin{table}[ht]
%\centering
%	\resizebox{\textwidth}{!}{ 
%\begin{tabular}{|l|c|c|c|c|c|c|}
%  \hline
%Модель Солоу & \multicolumn{2}{c|}{Первая группа} & \multicolumn{2}{c|}{Вторая группа} & \multicolumn{2}{c|}{Третья группа} \\ 
%с чел. капиталом & GDP PPP & GDP PC & GDP PPP & GDP PC & GDP PPP & GDP PC \\ 
%  \hline
%Оценка $\beta$ & 0.26 & 0.30 & 0.31 & 0.37 & 0.20 & 0.20 \\ 
%  Левая граница & 0.25 & 0.29 & 0.30 & 0.34 & 0.20 & 0.20 \\ 
%  Правая граница & 0.26 & 0.32 & 0.33 & 0.44 & 0.21 & 0.21 \\ 
%   \hline
%\end{tabular} }
%\end{table}
%	
%\end{frame}

\end{document}